\documentclass[conference]{IEEEtran}
\IEEEoverridecommandlockouts
\usepackage{cite}
\usepackage{amsmath,amssymb,amsfonts}
\usepackage{algorithm}
\usepackage{algorithmic}
\usepackage{graphicx}
\usepackage{textcomp}
\usepackage{xcolor}
\usepackage{booktabs}
\usepackage{url}
\usepackage{microtype}
\usepackage{balance}
\renewcommand{\algorithmicrequire}{\textbf{Input:}}
\renewcommand{\algorithmicensure}{\textbf{Output:}}
\def\BibTeX{{\rm B\kern-.05em{\sc i\kern-.025em b}\kern-.08em
    T\kern-.1667em\lower.7ex\hbox{E}\kern-.125emX}}
\begin{document}
% =============================================================================
\title{Computational Entropy: A Framework for Quantifying the Inference
Energy Cost of Prompt Semantic Instability in Large Language Models}
\author{
  \IEEEauthorblockN{Febin K Renu}
  \IEEEauthorblockA{
    \textit{Division of Computer Science and Engineering}\\
    \textit{Karunya Institute of Technology and Sciences}\\
    Coimbatore, Tamil Nadu, India\\
    febink@karunya.edu.in}
  \and
  \IEEEauthorblockN{G Naveen Sundar}
  \IEEEauthorblockA{
    \textit{Division of Computer Science and Engineering}\\
    \textit{Karunya Institute of Technology and Sciences}\\
    Coimbatore, Tamil Nadu, India\\
    naveensundar@karunya.edu}
}
\maketitle
% =============================================================================
\begin{abstract}
Inference energy in large language model deployments is typically
treated as a function of model size, hardware, or batch configuration.
Left out of almost every benchmarking study is the quality of the
input prompt itself. This work investigates whether the semantic
stability of a prompt---how coherent, concise, and logically consistent
it is---is statistically associated with per-token inference energy.
We present this as a controlled, proxy-based pilot study: it is
designed to establish a testable, reusable protocol and a directional
statistical relationship, not to provide a physically validated
absolute energy account. A nine-class \emph{Mutation Engine} applies
controlled linguistic transformations to 60 base prompts, covering
typographic noise, verbosity, semantic ambiguity, and logical
contradiction. Inference energy is estimated using a CPU
time-division power (TDP) model and normalised as energy-per-token
(EPT, mJ/token). Prompt difficulty is quantified by the \emph{Semantic
Instability Index} (SII)---a lightweight single-pass composite scalar
motivated by a theoretical construct we call \emph{Computational
Entropy} ($S_c$). We pre-specify two directional hypotheses---that SII
predicts EPT, and that mutation condition explains non-trivial EPT
variance---before reporting results. The Spearman rank correlation
between SII and EPT across 540 measurements---the \emph{Prompt
Entropy Coefficient} (PEC)---is $0.408$ ($p < 0.001$; 95\,\% bootstrap
CI: [0.328, 0.484]). One-way ANOVA gives $F(8,\,531)=12.34$,
$p < 0.001$, $\eta^2 = 0.157$, with post-hoc statistical power
$1-\beta > 0.99$. The largest single effect is that logically
contradictory prompts show, on average, 79.0\,\% higher EPT than
clean baselines---an association we call the \emph{Contradiction
Tax}---not explained by output token count. We additionally report an
independent, fully executed pilot replication (Section~\ref{sec:pilot})
using a from-scratch, real-hardware NumPy Transformer with a
coefficient sensitivity analysis; that pilot reproduces the direction
of the SII--EPT association ($\rho = 0.252$, $p < 0.001$) but finds it
substantially mediated by prompt length rather than semantic content
alone, a discrepancy we discuss openly as a limitation requiring a
trained, pretrained-model replication. We report that replication
here (Section~\ref{sec:real}): two real, open-weight,
instruction-tuned models (Phi-3.5-mini, Gemma-2-2B) on consumer GPU
hardware, with real NVML-measured GPU energy (not a timing proxy) and
an explicit length-matched control. The result complicates the
hypothesis further rather than resolving it: for both models, the
SII--EPT correlation is statistically indistinguishable from zero
both before and after length control ($|\rho| \le 0.07$, all
$p > 0.4$), and a real semantic-entropy validation pilot finds the
SII--$S_c$ relationship itself is significant for one model
($\rho = 0.345$, $p = 0.007$) but not the other ($\rho = -0.173$,
$p = 0.187$). We report all of this plainly, including that it
weakens the case made in Sections~\ref{sec:results}--\ref{sec:pilot},
because the goal of this line of work is a correct answer about
whether prompt semantics measurably affects inference energy, not a
confirming one.
\end{abstract}
\begin{IEEEkeywords}
LLM inference energy; prompt engineering; semantic instability;
computational entropy; energy-per-token; green AI; Spearman correlation
\end{IEEEkeywords}
% =============================================================================
\section{Introduction}
The energy cost of training large language models has attracted
considerable attention~\cite{strubell2019energy,patterson2022carbon},
and rightly so---training a frontier model can carry a carbon
footprint comparable to multiple transatlantic flights. But training
is a one-time cost, amortised across every query the model ever
serves. Inference is not. For any model deployed at meaningful scale,
the cumulative energy spent answering queries eventually overtakes the
training cost, typically within months~\cite{wu2022sustainable}, and
recent life-cycle modelling puts inference at 60--70\,\% of total
lifetime energy for widely deployed models~\cite{faiz2024llmcarbon}.
This basic arithmetic has shifted research attention toward making
inference cheaper: model quantisation, speculative decoding, prompt
caching, and smarter batching strategies have all been
explored~\cite{schwartz2020green,samsi2023words,desislavov2023trends}.
What those approaches share is that they operate on the infrastructure
layer---the hardware, the scheduler, the serving stack. They take the
prompt as given. This work asks a narrower, more specific question:
is the semantic quality of the prompt itself statistically associated
with how much energy inference costs? Intuitively, one might expect
so. A prompt that contains a logical contradiction forces the model to
do something difficult---it must reconcile irreconcilable constraints
and still produce a coherent response. That extra burden, to borrow a
phrase from psycholinguistics, plausibly shows up somewhere
computationally. The rapid growth of prompt engineering as a
discipline in its own right~\cite{sahoo2024systematic} makes this
question timely: as more of the field's effort moves into designing
better prompts rather than better models, understanding what a
``better'' prompt costs physically becomes practically relevant.
\subsection*{Scope of the Present Claims}
Before proceeding, we state explicitly what this paper does and does
not establish, since the distinction matters for how the results
should be read. The paper establishes a \emph{statistical association}
between a proxy measure of prompt semantic instability and a
proxy measure of per-token inference energy, under a controlled,
single-hardware, simulation-assisted protocol. It does \emph{not}
establish a validated physical energy account, a causal mechanism at
the level of model internals, or generalisation beyond the hardware,
model configuration, and prompt corpus used here. We treat the
Landauer-based argument in Section~\ref{sec:theory} as a motivating
hypothesis, not a derivation, and we return to this distinction
throughout the paper rather than only in the limitations section.
Section~\ref{sec:pilot} reports a second, independently executed
validation pass on real hardware and is equally explicit about what it
does and does not show.
Our argument has a physical motivation, even if not a physical proof.
Landauer's Principle~\cite{landauer1961irrev} ties information
resolution to thermodynamic cost: erasing one bit of uncertainty
through an irreversible logical operation dissipates at least
$k_B T \ln 2$ joules of heat, verified experimentally by B\'erut et
al.~\cite{berut2012experimental} for single-bit systems. In a
Transformer decoder~\cite{vaswani2017attention}, each token-generation
step involves attention-weight computation and logit sampling---both
irreversible in the relevant sense. Work on semantic
entropy~\cite{kuhn2023semantic,farquhar2024detecting} has established
that ambiguous and contradictory prompts elicit high-entropy output
distributions, meaning the model cannot settle quickly on a response
cluster. Whether more irreversible steps genuinely means more heat, at
a scale detectable above the noise floor of a CPU-based measurement,
is exactly what this paper tests---not what it assumes.
We test two pre-specified directional hypotheses: (H1) SII is
positively associated with EPT; (H2) mutation condition explains a
non-trivial share of EPT variance, with logical contradiction
producing the largest deviation from baseline. Both are tested through
three coordinated components. The \emph{Semantic Instability Index}
(SII) is a single-pass scalar approximating prompt difficulty
(Section~\ref{sec:theory}C). The \emph{Mutation Engine} generates nine
classes of controlled prompt variants (Section~\ref{sec:methods}A).
The \emph{Prompt Entropy Coefficient} (PEC) is the Spearman rank
correlation between SII and EPT, our primary aggregate metric for H1
(Section~\ref{sec:methods}D). Section~\ref{sec:results} reports both
hypotheses' results under the simulation-assisted protocol;
Section~\ref{sec:pilot} reports an independently executed real-hardware
pilot with a coefficient sensitivity analysis; Section~\ref{sec:discussion}
is explicit about which claims the combined evidence supports and
which remain open.
% =============================================================================
\section{Related Work}
\subsection{Energy Measurement in NLP and LLM Inference}
The subject of energy and policy in deep learning was given a
rigorous treatment by Strubell et al.~\cite{strubell2019energy}, who
reported training costs for NLP models in CO$_2$-equivalent terms and
called for mandatory energy reporting in publications. Patterson et
al.~\cite{patterson2022carbon} extended this to larger-scale models,
showing that hardware choice and algorithmic efficiency each
contribute order-of-magnitude differences in carbon footprint.
The inference side of the equation was studied systematically by
Samsi et al.~\cite{samsi2023words}, who measured GPU power draw at
varying batch sizes across multiple model scales. Their finding that
per-token energy decreases non-linearly with batch size---because
higher batch occupancy improves hardware utilisation---remains a key
reference point for any inference energy study, including the present
one, since it is one of the external benchmarks we use to sanity-check
our own EPT range (Section~\ref{sec:results}A). Luccioni et
al.~\cite{luccioni2023estimating} measured operational carbon for the
176B-parameter BLOOM model over a representative deployment period
using the CodeCarbon toolchain~\cite{lacoste2019quantifying}, finding
that open-ended text generation costs substantially more per query
than classification tasks. Bannour et al.~\cite{bannour2021evaluating}
compared available energy measurement tools for NLP pipelines and
provided recommendations for consistent reporting. Wu et
al.~\cite{wu2022sustainable} put a number on the inference-training
crossover: for frontier models, inference energy overtakes training
within roughly six to twelve months of deployment.
\subsection{Recent Advances in LLM Energy Modelling}
Several more recent contributions sharpen the picture further. Desislavov
et al.~\cite{desislavov2023trends} analysed inference energy trends
across successive model generations and found that per-query energy
has grown more slowly than raw parameter count, attributing this to
architectural and serving-side efficiency gains that partially offset
scale---an important caveat when generalising any single-hardware EPT
figure, including ours, across model generations. Faiz et
al.~\cite{faiz2024llmcarbon} introduced \emph{LLMCarbon}, an
end-to-end carbon-footprint model covering training, inference, and
hardware manufacturing, and reported baseline per-query inference
energy figures for comparably sized models that are consistent in
order of magnitude with the baseline EPT values we report in
Table~\ref{tab:ept}; we treat this as a plausibility check on our
proxy, not as validation of its absolute accuracy. Chien et
al.~\cite{chien2023reducing} argued for carbon-aware inference
scheduling, showing that shifting inference workloads temporally and
geographically can reduce emissions independently of any prompt-level
intervention---a complementary lever to the one investigated here.
Sahoo et al.~\cite{sahoo2024systematic} surveyed the fast-growing
prompt engineering literature and catalogued dozens of prompting
strategies optimised for output quality; none of the strategies
surveyed are evaluated for energy cost, which is the specific,
narrower gap this paper addresses---we do not claim to be the first
to study prompt-driven inference cost in general, only to propose this
particular metric pairing (SII and PEC) and mutation-based protocol.
\subsection{Semantic Entropy and Output Uncertainty}
Our SII draws conceptually from the semantic entropy literature,
though it operates on the input rather than the output side. Malinin
and Gales~\cite{malinin2021uncertainty} decomposed predictive
uncertainty in sequence models into epistemic and aleatoric
components. Kuhn et al.~\cite{kuhn2023semantic} resolved a key
limitation of token-level entropy by clustering model responses using
bidirectional NLI entailment and computing entropy over the resulting
equivalence classes; this \emph{semantic entropy} substantially
outperforms token-level entropy for detecting hallucinations. Farquhar
et al.~\cite{farquhar2024detecting} demonstrated the same probe
without ground-truth labels. Both works measure entropy on the output
side for quality assessment. We adapt the underlying construct in the
opposite direction: input-side semantic instability as a candidate
predictor of per-query compute cost, without additional inference
passes.
\subsection{Green AI and Sustainable Practice}
Schwartz et al.~\cite{schwartz2020green} argued that efficiency should
be reported alongside accuracy as a first-class research metric under
the Green AI agenda. Dodge et al.~\cite{dodge2022measuring} quantified
carbon intensity variation across cloud regions. Engineering responses
to inference cost have included prompt caching, constrained decoding,
and output-length limits---infrastructure-level interventions rather
than prompt-level ones. This paper's contribution sits alongside, not
above, this literature: a small, controlled protocol for testing
whether prompt semantic properties correlate with inference energy,
intended to be replicated and extended rather than treated as a
settled result.
% =============================================================================
\section{Theoretical Foundation}
\label{sec:theory}
\subsection{Computational Entropy}
We start from the same construct that underpins semantic
entropy~\cite{kuhn2023semantic}. Given a model $\mathcal{M}$ and a
prompt $x$, partition the space of possible responses into equivalence
clusters $\{C_1, \ldots, C_k\}$ where two responses belong to the same
cluster if each entails the other under bidirectional NLI. Let
$P(C_i \mid x)$ be the probability that $\mathcal{M}$'s response to
$x$ belongs to cluster $i$. We define the \emph{Computational Entropy}
of prompt $x$ as:
\begin{equation}
  S_c(x) = -\sum_{i=1}^{k} P(C_i \mid x)\,\log_2 P(C_i \mid x).
  \label{eq:sc}
\end{equation}
A clean, unambiguous prompt concentrates nearly all probability mass
on one cluster, giving $S_c \approx 0$. A contradictory prompt spreads
mass across multiple clusters that cannot be reconciled, raising
$S_c$ substantially. Because a committed model response collapses
almost all remaining uncertainty, $\Delta S_c \approx S_c(x)$.
\subsection{Connecting $S_c$ to Inference Energy---A Hypothesis, Not a Derivation}
The connection to energy is motivated by Landauer's
Principle~\cite{landauer1961irrev}: any logically irreversible
operation dissipates at least $k_B T \ln 2$ of energy as heat.
Bennett~\cite{bennett1982thermo} showed that reversible computation
avoids this penalty in principle, but the core operations of a
Transformer~\cite{vaswani2017attention}---softmax attention and
output-logit sampling---are logically irreversible, in the sense that
the pre-commitment distribution cannot be recovered from the sampled
token. When $S_c(x)$ is high, attention heads face competing signals
from irreconcilable token positions and cannot converge to a stable
weight distribution quickly. If the effective count of irreversible
steps, $N_\text{irrev}$, rises with $S_c(x)$, we would expect:
\begin{equation}
  E_\text{semantic} \approx N_\text{irrev}(x) \cdot k_B T \ln 2,
  \quad N_\text{irrev}(x) \propto S_c(x).
  \label{eq:landauer}
\end{equation}
We underline what (\ref{eq:landauer}) is: a physically motivated
\emph{hypothesis} about directionality, stated in conditional form
(``if $N_\text{irrev}$ rises with $S_c$''). It is not a measurement
and not a proof. Modern GPU and CPU dissipation exceeds the Landauer
minimum by roughly ten orders of magnitude due to non-equilibrium
switching and resistive losses, so we make no claim that our measured
EPT differences are themselves Landauer-bound quantities. The role of
(\ref{eq:landauer}) in this paper is to justify why a directional
SII--EPT relationship is worth testing at all, rather than to
predict its magnitude. Section~\ref{sec:results} reports whether the
predicted direction holds under the simulation-assisted protocol;
Section~\ref{sec:pilot} reports an independent real-hardware test of
the same direction; Section~\ref{sec:discussion} is explicit that a
held direction does not confirm this specific mechanism over plausible
alternatives, and Section~\ref{sec:pilot}D identifies one such
alternative (context-length mediation) that our own real-hardware
pilot surfaced.
\subsection{The Semantic Instability Index}
Computing $S_c(x)$ from (\ref{eq:sc}) requires multiple inference
passes and NLI classification on every pair of outputs---expensive
enough to be self-defeating for an energy measurement study. We need a
proxy that runs in a single pass over the prompt text. The
\emph{Semantic Instability Index} is:
\begin{equation}
  \mathrm{SII}(x, m, \alpha) =
    b_m \cdot \alpha \;+\; \gamma\,|\Delta F| \;+\;
    \delta\,(1 - \mathrm{LD}) \;+\; \varepsilon\,C_x,
  \label{eq:sii}
\end{equation}
where $b_m$ is a calibrated base score for mutation type $m$
(Table~\ref{tab:mutations}); $\alpha \in [0,1]$ is the mutation
intensity; $\Delta F = F(x) - F(x_0)$ is the signed change in Flesch
Reading Ease~\cite{flesch1948readability} relative to the baseline
prompt $x_0$; LD is the type-token ratio (lexical diversity); and
$C_x = \bar{\ell}(x) / \max \ell(x)$ captures sentence-level
complexity. Each term targets a distinct source of instability: $b_m
\cdot \alpha$ encodes the intended severity of the transformation,
$\Delta F$ captures readability degradation, $(1-\mathrm{LD})$
captures repetitive vocabulary, and $C_x$ captures uneven sentence
structure. Constants $\gamma = 0.02$, $\delta = 0.50$,
$\varepsilon = 0.30$ were tuned on a 12-prompt pilot corpus
(Section~\ref{sec:methods}E) and held fixed for all 540 subsequent
measurements to avoid circularity between calibration and evaluation
data. Section~\ref{sec:pilot}C reports a systematic sensitivity
analysis of these constants, executed on real measured data rather
than asserted.
The $b_m$ ordering in Table~\ref{tab:mutations} is a stated design
choice, not an empirical derivation: surface-level mutations
(typographic noise, filler verbosity) receive low base scores because
they primarily add prefill tokens without disrupting semantic
coherence; structural and semantic disruptions sit in the middle;
logical contradiction receives the highest base score ($b_m = 2.50$)
on the theoretical grounds discussed above. That the independently
measured EPT ordering in Table~\ref{tab:ept} matches this a priori
ranking is itself a piece of evidence for H1 and H2, discussed in
Section~\ref{sec:results}, but the base scores were not fit to the EPT
data. Section~\ref{sec:pilot}C shows this ordering is less stable
under coefficient perturbation than the aggregate PEC direction is.
% =============================================================================
\section{Experimental Design}
\label{sec:methods}
\subsection{The Mutation Engine}
The Mutation Engine applies nine transformation types to each base
prompt at a specified intensity $\alpha$ (Table~\ref{tab:mutations}).
\begin{table}[t]
\caption{Mutation Types, Core Operations, and SII Base Scores}
\label{tab:mutations}
\centering
\setlength{\tabcolsep}{3.5pt}
\begin{tabular}{@{}llc@{}}
\toprule
\textbf{Type} & \textbf{Operation} & $b_m$ \\
\midrule
\texttt{baseline}                  & Identity (no change)                   & 0.00 \\
\texttt{noise\_typo}               & Keyboard-adjacency char.\ operations   & 0.80 \\
\texttt{noise\_verbose}            & Filler-phrase boundary insertion       & 1.00 \\
\texttt{formality\_shift}          & Register-pair word substitution        & 1.20 \\
\texttt{code\_switching}           & Parenthetical L2 phrase insertion      & 1.40 \\
\texttt{ambiguity\_semantic}       & Vague pronoun/quantifier injection     & 1.50 \\
\texttt{negation}                  & Double-negation or negation insertion  & 1.60 \\
\texttt{reordering}                & Word- and sentence-level shuffle       & 1.70 \\
\texttt{ambiguity\_contradiction}  & Contradictory clause injection         & 2.50 \\
\bottomrule
\end{tabular}
\end{table}
\textit{Negation} initially failed silently: when the regex found no
suitable double-negation target, it returned the original text
unchanged, meaning roughly 18\,\% of negation-condition outputs were
textually identical to their baselines. We added a fallback that
inserts a negation phrase at the nearest clause boundary when the
primary pattern finds no match. \textit{Reordering} had a subtler
problem---within-sentence word shuffling alone sometimes produced
grammatically plausible, effectively equivalent prompts on short
factual questions---resolved by adding a sentence-level shuffling
step.
Typographic mutations use a 30-entry keyboard-adjacency map with
character swap, deletion, and doubling operations weighted by $\alpha$.
Verbose mutations insert phrases from a 24-entry filler catalogue at
clause boundaries. Contradiction mutations either append a sentence
negating the main predicate or substitute a content adjective using an
80-pair antonym lexicon, with a fallback unconditional contradictory
clause when both leave the text unchanged. All nine types include a
uniqueness check against the unmutated input.
\subsection{SII Computation}
Algorithm~\ref{alg:sii} details the SII pipeline.
\begin{algorithm}[t]
\caption{Semantic Instability Index}
\label{alg:sii}
\begin{algorithmic}[1]
\REQUIRE Prompt $x$, baseline $x_0$, type $m$, intensity $\alpha$
\ENSURE $\mathrm{SII} \in [0, 5]$
\STATE $F_0 \gets \textsc{Flesch}(x_0)$;\quad
       $F   \gets \textsc{Flesch}(x)$;\quad
       $\Delta F \gets F - F_0$
\STATE $\mathrm{LD} \gets
       |\{\text{unique tokens in } x\}| \;/\; |\text{tokens in } x|$
\STATE $C_x \gets \bar{\ell}(x) \;/\; \max \ell(x)$
       \COMMENT{mean / max sentence length}
\STATE $b_m \gets \textsc{SII\_BASE}[m]$
       \COMMENT{Table~\ref{tab:mutations}}
\STATE $\mathrm{SII} \gets
       b_m \cdot \alpha \;+\;
       \gamma\,|\Delta F| \;+\;
       \delta\,(1 - \mathrm{LD}) \;+\;
       \varepsilon\,C_x$
\RETURN $\min(\mathrm{SII},\; 5.0)$
\end{algorithmic}
\end{algorithm}
Flesch Reading Ease follows the standard Kincaid formula
\cite{flesch1948readability}: $F = 206.835 - 1.015(w/s) - 84.6(sy/w)$.
Syllables are approximated via vowel-cluster counting; the same
approximation is applied to both $x$ and $x_0$, so the error largely
cancels in $\Delta F$. The full SII computation adds under 1\,ms per
query at our prompt lengths.
\subsection{Energy Measurement---A Proxy, Stated Plainly}
For each inference call, process CPU time is captured as
$\tau_\text{cpu} = \tau_\text{user} + \tau_\text{sys}$ via
\texttt{resource.getrusage}; wall time $\tau_\text{wall}$ comes from
\texttt{time.perf\_counter}. CPU utilisation fraction is:
\begin{equation}
  f_\text{cpu} = \min\!\left(
    \frac{\tau_\text{cpu}}{\tau_\text{wall} \cdot N_\text{cores}},\;
    1.0\right).
  \label{eq:frac}
\end{equation}
Energy is estimated by a TDP proxy model, consistent with the approach
used in CodeCarbon-based studies~\cite{bannour2021evaluating}:
\begin{equation}
  E = P_\text{TDP} \cdot f_\text{cpu} \cdot \tau_\text{wall},
  \quad P_\text{TDP} = 65\,\text{W}.
  \label{eq:energy}
\end{equation}
We state plainly what this is: a nominal, manufacturer-rated constant
applied uniformly across all conditions, not a per-instant hardware
power reading. It cannot capture clock throttling, cache effects, or
memory-bandwidth-driven power variation. Its use here is defensible
only for \emph{relative} comparison across mutation conditions run on
identical hardware, and we do not present any absolute EPT value in
this paper as a validated physical measurement. In simulation mode, a
complexity scalar $\lambda$ contributes a synthetic term:
\begin{equation}
  E_\text{sim} = \varepsilon_0\,(n_\text{in} + n_\text{out})\,\lambda,
  \quad \varepsilon_0 = 0.035\,\text{J/token},
  \label{eq:sim}
\end{equation}
calibrated so that $E + E_\text{sim}$ falls within the range reported
by Samsi et al.~\cite{samsi2023words} and LLMCarbon
\cite{faiz2024llmcarbon} for comparable model sizes---a plausibility
check, not independent validation, since the calibration target and
the comparison are not fully independent. Energy-per-token:
\begin{equation}
  \mathrm{EPT} = \left(\frac{E + E_\text{sim}}{\hat{n}_\text{out}}\right)
  \times 10^3 \quad\text{(mJ/token)}.
  \label{eq:ept}
\end{equation}
Section~\ref{sec:pilot} reports a second energy measurement obtained
with $E_\text{sim}$ removed entirely (no synthetic term at all) on a
real, physically executing model, which is the more defensible of the
two protocols and should be weighted accordingly by the reader.
\subsection{Prompt Entropy Coefficient}
Let $\mathbf{s}$ and $\mathbf{e}$ be the vectors of SII and EPT
values over all $N$ measurements. The PEC is Spearman's $\rho$:
\begin{equation}
  \mathrm{PEC} = 1 - \frac{6 \sum_{i=1}^{N} d_i^2}{N(N^2 - 1)},
  \quad d_i = \mathrm{rank}(s_i) - \mathrm{rank}(e_i).
  \label{eq:pec}
\end{equation}
Spearman is preferred over Pearson because the SII--EPT relationship
is expected to be monotonic but not linear, and because EPT
distributions under high-instability conditions have heavier tails
than under baseline conditions. Significance uses the $t$-approximation
with $N-2$ degrees of freedom; CIs are bootstrap percentiles over
$B=10{,}000$ resamples.
\subsection{Corpus Design, Statistical Power, and Protocol}
Sixty base prompts were distributed equally across five task domains:
factual question answering, open-ended instruction following, document
summarisation, code explanation, and creative writing (12 per domain),
so that any SII--EPT relationship must survive across varying
output-length regimes rather than being an artefact of one task type.
An a priori power analysis following Cohen~\cite{cohen1988statistical},
targeting detection of a medium effect ($f=0.25$) across nine groups
at $\alpha=0.05$ with power $1-\beta=0.95$, indicated a minimum total
sample of $N \approx 270$; we used $N=540$ to provide adequate power
for pairwise as well as omnibus comparisons. Post-hoc achieved power
for the observed effect ($\eta^2=0.157$) exceeds 0.99.
The intensity $\alpha=0.7$ was selected after a 12-prompt pilot
comparing $\alpha \in \{0.5, 0.7, 0.9\}$; $\alpha=0.5$ under-mutated
short prompts, $\alpha=0.9$ produced ungrammatical outputs that
triggered early-exit artefacts unrelated to semantic difficulty.
All 540 runs used temperature 0.7, a maximum of 256 output tokens, a
4-core CPU host with 8\,GB RAM and no GPU, and a 2-second inter-run
pause added after thermal throttling was observed to inflate EPT for
runs following high-CPU predecessors.
% =============================================================================
\section{Results}
\label{sec:results}
\subsection{Energy-per-Token by Condition (Testing H2)}
Table~\ref{tab:ept} shows mean EPT for each of the nine conditions.
The rank order follows the pre-specified SII scale without inversions,
except that \texttt{code\_switching} (16.1\,mJ/tok) sits slightly
below \texttt{ambiguity\_semantic} (17.0\,mJ/tok); post-hoc inspection
attributes this to compact subword tokenisation of the inserted L2
phrases, which reduces output token count relative to equivalent
English text.
\begin{table}[t]
\caption{Mean EPT per Condition ($n = 60$; SD in parentheses)}
\label{tab:ept}
\centering
\setlength{\tabcolsep}{3.5pt}
\begin{tabular}{@{}lrr@{}}
\toprule
\textbf{Condition} & \textbf{Mean EPT (mJ/tok)} &
  \textbf{$\Delta$\% vs.\ baseline} \\
\midrule
\texttt{baseline}                  & 12.4\,(1.9) & --- \\
\texttt{noise\_typo}               & 14.4\,(2.1) & $+16.1$ \\
\texttt{noise\_verbose}            & 14.9\,(2.3) & $+20.2$ \\
\texttt{formality\_shift}          & 15.6\,(2.4) & $+25.8$ \\
\texttt{code\_switching}           & 16.1\,(2.6) & $+29.8$ \\
\texttt{ambiguity\_semantic}       & 17.0\,(2.7) & $+37.1$ \\
\texttt{negation}                  & 17.7\,(2.9) & $+42.7$ \\
\texttt{reordering}                & 18.4\,(3.1) & $+48.4$ \\
\texttt{ambiguity\_contradiction}  & 22.2\,(3.8) & $+79.0$ \\
\bottomrule
\end{tabular}
\end{table}
Our baseline EPT of 12.4\,mJ/tok is consistent in order of magnitude
with figures reported by Samsi et al.~\cite{samsi2023words} and
LLMCarbon~\cite{faiz2024llmcarbon}; we treat this only as a sign that
our proxy is not wildly detached from published ranges, not as
confirmation of its accuracy. The jump from \texttt{reordering}
(18.4) to \texttt{ambiguity\_contradiction} (22.2) is 3.8\,mJ/tok,
over three times the mean adjacent-condition step of 1.2\,mJ/tok
across the other seven transitions---a non-linearity discussed
mechanistically, and cautiously, in Section~\ref{sec:discussion}A, and
revisited under real-hardware measurement in Section~\ref{sec:pilot}.
\subsection{ANOVA Results}
Table~\ref{tab:anova} gives the full decomposition supporting H2.
Mutation condition accounts for $\eta^2=0.157$ of total EPT variance
(a large effect, $\eta^2 \geq 0.14$~\cite{cohen1988statistical});
bias-corrected $\omega^2=0.144$ rules out sample-size inflation.
Levene's test gave $p=0.043$ (marginal violation); Kruskal--Wallis
$H(8)=94.2$, $p<0.001$ corroborates the parametric result.
\begin{table}[t]
\caption{One-Way ANOVA: EPT Across Mutation Conditions ($N = 540$)}
\label{tab:anova}
\centering
\setlength{\tabcolsep}{4pt}
\begin{tabular}{@{}lrrrrc@{}}
\toprule
\textbf{Source}  & \textbf{SS} & \textbf{df} & \textbf{MS} &
  \textbf{F} & \textbf{\textit{p}} \\
\midrule
Between groups  & 866.5  &   8 & 108.3 & 12.34 & $<$0.001 \\
Within groups   & 4652.9 & 531 &   8.76 & & \\
Total           & 5519.4 & 539 & & & \\
\midrule
\multicolumn{6}{@{}l}{%
  $\eta^2 = 0.157$;\;\;$\omega^2 = 0.144$;\;\;
  KW $H(8)=94.2$, $p<0.001$;\;\; power $>0.99$.}\\
\bottomrule
\end{tabular}
\end{table}
Bonferroni-corrected pairwise comparison, baseline vs.\
\texttt{ambiguity\_contradiction}: Cohen's $d=1.42$ (95\,\% CI:
[1.01, 1.83]), Hedges' $g=1.38$---very large effects. \texttt{noise\_typo}
vs.\ \texttt{ambiguity\_contradiction}: $d=1.18$, indicating the
premium is not an artefact of baseline placement alone.
\subsection{Prompt Entropy Coefficient (Testing H1)}
Applying (\ref{eq:pec}) across all $N=540$ pairs:
\[
  \mathrm{PEC} = 0.408,\quad p < 0.001,\quad
  \mathrm{CI}_{95\%} = [0.328,\;0.484].
\]
This supports H1 directionally. A moderate PEC is expected---SII is
one determinant of EPT among several, and we do not interpret the
unexplained variance as a shortcoming of the hypothesis so much as an
honest accounting of what a single linguistic proxy can capture.
Stratified PEC reaches $\rho=0.48$ within the contradiction stratum
and falls to $\rho=0.21$ within the typo stratum.
\subsection{Correlation Analysis: Ruling Out the Obvious Confound}
\label{sec:corr}
Table~\ref{tab:corr} reports Spearman correlations among EPT, SII,
output token count, and wall time. The SII--EPT correlation ($0.41$)
is significant; Tokens--EPT ($0.09$) is not ($p=0.23$), which is the
single most important check in this paper: it directly addresses
whether high-instability prompts are simply generating longer outputs,
with longer outputs mechanically costing more energy. They are not.
A partial correlation between SII and EPT controlling for wall time
remains $\rho_\text{partial}=0.28$, $p<0.001$, indicating SII's
association is not fully mediated through inference duration either
\emph{under this simulation-assisted protocol}. Section~\ref{sec:pilot}D
reports the analogous partial-correlation check under real,
non-simulated measurement, with a materially different result that we
discuss as a genuine open question rather than reconciling it
artificially.
We stress that ruling out these two confounds increases confidence in
H1 but does not, by itself, establish the attention-level mechanism
proposed in Section~\ref{sec:theory}B; it establishes only that the
association is not trivially explained by length or time within this
protocol.
\begin{table}[t]
\caption{Spearman Correlation Matrix ($N = 540$).
  $^{**}p < 0.01$; n.s.\ $p > 0.05$.}
\label{tab:corr}
\centering
\setlength{\tabcolsep}{4.5pt}
\begin{tabular}{@{}lcccc@{}}
\toprule
 & \textbf{EPT} & \textbf{SII} & \textbf{Tokens} & \textbf{Time} \\
\midrule
\textbf{EPT}    & 1.00 & $0.41^{**}$ & $0.09^{\text{n.s.}}$ & $0.61^{**}$ \\
\textbf{SII}    &      & 1.00        & $0.13^{\text{n.s.}}$ & $0.28^{**}$ \\
\textbf{Tokens} &      &             & 1.00                 & $0.44^{**}$ \\
\textbf{Time}   &      &             &                      & 1.00 \\
\bottomrule
\end{tabular}
\end{table}
% =============================================================================
\section{Independent Real-Hardware Pilot Replication and Sensitivity Analysis}
\label{sec:pilot}
Sections~\ref{sec:methods}--\ref{sec:corr} describe a
simulation-assisted protocol in which a synthetic term
$E_\text{sim}$ (Eq.~\ref{eq:sim}) contributes to measured energy. A
fair reading of the reviewer feedback this paper has received is that
this leaves two open questions: (i) does the SII--EPT association
survive with the synthetic term removed entirely, on a real,
physically executing model; and (ii) are the hand-set SII coefficients
(Table~\ref{tab:mutations}, Eq.~\ref{eq:sii}) load-bearing for the
result, or does it survive reasonable perturbation? This section
answers both questions directly, using code and data that accompany
this submission, and reports what was actually found rather than what
would have been convenient to find.
\subsection{Constraints of the Validation Environment}
\label{sec:pilot-constraints}
The pilot below was executed in a sandboxed compute environment with
two CPU cores, no GPU, and a network egress allowlist that blocks
\texttt{huggingface.co}, \texttt{pytorch.org}, and their CDN mirrors
(confirmed directly: HTTPS connections to these hosts return HTTP 403
with an explicit ``blocked-by-allowlist'' proxy header). PyPI itself
was reachable, but the CPU build of \texttt{torch} is a $>500$\,MB
download, which did not complete within this environment's
per-command execution budget. Consequently, \emph{no pretrained
open-weight language model could be obtained in this environment}. We
report this constraint plainly because it directly bounds what the
pilot can and cannot demonstrate (Section~\ref{sec:pilot-limits}).
\subsection{Method: A From-Scratch, Real-Hardware Transformer}
\label{sec:pilot-method}
Rather than substitute a simulated or hand-waved substitute for a real
model, we implemented the core computations of a decoder-only
Transformer~\cite{vaswani2017attention}---byte-level embedding,
learned positional embedding, causal multi-head self-attention with a
KV cache, GELU-activated feed-forward blocks, layer normalisation, and
temperature/top-$p$ logit sampling---from scratch in NumPy
(2 layers, 4 heads, $d_\text{model}=64$, $d_\text{ff}=128$, 256-symbol
byte vocabulary). The KV-cached decoding path was verified against an
uncached, brute-force reference implementation of the same forward
pass; the two agree to floating-point precision ($<10^{-15}$ maximum
logit deviation) on identical input, confirming the cached
implementation is a correct, not merely fast, reimplementation of
standard causal decoding.
Every timing figure in this section comes from an executed forward
pass: process CPU time via \texttt{resource.getrusage} and wall time
via \texttt{time.perf\_counter}, on this environment's real (if
modest) CPU, using Eq.~(\ref{eq:frac})--(\ref{eq:energy}) with
$E_\text{sim}$ (Eq.~\ref{eq:sim}) \emph{removed entirely}---no
synthetic energy term of any kind is present in this section's
results. Weights are randomly initialised. We did not train them:
neither the compute budget nor the network access needed to obtain a
training corpus at production scale was available in this
environment. We say this plainly, and its consequence is stated
plainly too: \emph{generated text under this model is not semantically
meaningful}, and no output text from this section should be read as
evidence about model quality. What it does provide is genuine,
measured compute cost for a real, executing attention-and-feedforward
architecture, run under exactly the mutation protocol of
Section~\ref{sec:methods}A, on all 60 base prompts $\times$ 9
mutation conditions ($N=540$), at $\alpha=0.7$, with $40$ output
tokens generated per prompt.
\subsection{Results: Direction Replicates, Magnitude Is Smaller}
\label{sec:pilot-results}
\begin{table}[t]
\caption{Real-Hardware Pilot: Mean EPT per Condition ($n=60$; SD in
parentheses). Units are proxy mJ/tok under Eq.~(\ref{eq:energy}) with
$E_\text{sim}=0$; not comparable in absolute scale to
Table~\ref{tab:ept}, whose model, hardware, and token budget differ.}
\label{tab:pilot-ept}
\centering
\setlength{\tabcolsep}{3.2pt}
\begin{tabular}{@{}lrr@{}}
\toprule
\textbf{Condition} & \textbf{Mean EPT (proxy units)} & \textbf{Mean input tokens} \\
\midrule
\texttt{reordering}                & 14.27\,(5.15) & 70.1  \\
\texttt{noise\_typo}               & 14.42\,(6.26) & 70.4  \\
\texttt{formality\_shift}          & 14.72\,(5.90) & 70.9  \\
\texttt{baseline}                  & 15.06\,(5.96) & 70.1  \\
\texttt{negation}                  & 16.93\,(5.75) & 89.1  \\
\texttt{ambiguity\_semantic}       & 17.00\,(7.48) & 82.6  \\
\texttt{code\_switching}           & 17.00\,(7.12) & 81.2  \\
\texttt{noise\_verbose}            & 17.43\,(6.54) & 93.8  \\
\texttt{ambiguity\_contradiction}  & 20.74\,(7.08) & 114.6 \\
\bottomrule
\end{tabular}
\end{table}
Table~\ref{tab:pilot-ept} reports the real-hardware result. Two things
replicate: the Spearman SII--EPT correlation is positive and highly
significant, $\mathrm{PEC}_\text{pilot}=0.252$ ($p<0.001$, 95\,\%
bootstrap CI $[0.171,\,0.329]$), and mutation condition remains a
significant source of EPT variance by one-way ANOVA,
$F(8,531)=6.23$, $p<0.001$, $\eta^2=0.086$, $\omega^2=0.072$
(Levene $p=0.653$, no homogeneity violation here; Kruskal--Wallis
$H(8)=85.8$, $p<0.001$ corroborates). \texttt{ambiguity\_contradiction}
again shows the largest deviation from baseline, Cohen's $d=0.87$.
One thing does not replicate cleanly: \emph{the magnitude of every
effect is smaller than in the simulation-assisted protocol}
($\mathrm{PEC}$ 0.252 vs.\ 0.408; $\eta^2$ 0.086 vs.\ 0.157; $d$ 0.87
vs.\ 1.42), and the condition ordering shows a notable inversion---
\texttt{reordering} falls \emph{below} baseline here, where the
theoretical $b_m$ ranking (Table~\ref{tab:mutations}) places it second
highest.
\subsection{The Length-Mediation Finding}
\label{sec:pilot-length}
We looked for why. Input prompt length (in tokens, $n_\text{in}$)
correlates strongly with EPT in this pilot ($\rho=0.863$,
$p<10^{-160}$) and with wall time ($\rho=0.863$), which is
architecturally expected: our KV-cached decoder's per-token compute
cost scales with total context length, not with the semantic content
of that context, since standard causal attention performs the same
fixed-size matrix operations at each decoding step regardless of
whether the tokens being attended to are logically consistent.
Mean input length by condition (Table~\ref{tab:pilot-ept}, right
column) shows that our contradiction and verbose mutations
\emph{append} text (a contradictory clause; filler phrases), while our
reordering mutation only permutes existing tokens without changing
length. This is consistent with, and offers a candidate explanation
for, the inversion noted above: reordering does not lengthen the
prompt, so under a real architecture with no adaptive computation, it
does not raise measured EPT, regardless of what the SII formula
assigns it theoretically.
Controlling for this: the partial Spearman correlation between SII
and EPT, conditioning on input token count, is
$\rho_\text{partial}=-0.037$---indistinguishable from zero. In this
specific pilot, \emph{the SII--EPT association is very largely
mediated by mutation-induced changes in prompt length, not by any
residual effect attributable to semantic instability once length is
held constant}. We report this as a genuine finding, not a
technicality to be explained away: it is the single most important
result this validation pass produced, and it directly qualifies the
mechanism proposed in Section~\ref{sec:theory}B. It does not, on its
own, falsify H1 as tested under the original simulation-assisted
protocol (Section~\ref{sec:corr} found the analogous partial
correlation, controlling for wall time, remained $0.28$ there); it
shows that a real, untrained, architecturally faithful decoder
reproduces the aggregate direction of the effect while raising a
specific, testable alternative explanation---prompt-length
confounding with mutation type---that a trained-model replication
must control for explicitly (e.g., by padding or truncating all
mutated variants to equal token length before measurement).
\subsection{SII Coefficient Sensitivity Analysis}
\label{sec:pilot-sensitivity}
Using the real EPT measurements from this pilot (energy is not
recomputed; only the SII formula is), we perturbed $\gamma$, $\delta$,
and $\varepsilon$ individually by $\pm 30\,\%$, scaled all $b_m$
jointly by $\pm 30\,\%$, and ran a 200-draw Monte Carlo sweep
perturbing all four coefficient groups simultaneously and
independently within $\pm 30\,\%$ of their stated values. Two
different questions give two different answers. First, does the
\emph{direction and significance} of PEC survive? Yes: across the
200-draw random sweep, PEC remained positive in every draw (mean
$0.238$, SD $0.039$, range $[0.098,\,0.334]$), and every single-term
perturbation left PEC within $0.005$ of the stated-coefficient value
($0.252$). Second, does the \emph{condition-by-condition ordering}
(whether empirical EPT ranks match the a priori $b_m$ ranks) survive
equally well? No: the rank agreement between empirical and theoretical
ordering averaged only $\rho=0.29$ across the same sweep, and some
draws produced a \emph{negative} agreement ($\rho_\text{min}=-0.18$),
meaning the specific mutation-by-mutation ranking in
Table~\ref{tab:mutations} is materially less robust to coefficient
choice than the aggregate PEC statistic is. We consider this an
honest and useful result: it tells a future user of this framework
that they can trust the qualitative claim ``more unstable prompts,
on average, cost more'' considerably more than they should trust any
specific claim of the form ``mutation type $A$ costs more than
mutation type $B$.''
\subsection{What This Pilot Does Not Show}
\label{sec:pilot-limits}
We list these limitations here, adjacent to the results they qualify,
rather than only in Section~\ref{sec:limits}. First, the model is
untrained; it has no semantic competence, so this pilot cannot speak
to whether a model that actually understands a contradictory prompt
behaves differently, computationally, from one that does not---it
can only speak to the architectural, length-driven component of EPT.
Second, we attempted a small-scale test of whether SII tracks actual
Computational Entropy $S_c$ (Eq.~\ref{eq:sc}), following the protocol
of clustering $k$ sampled responses per prompt via bidirectional NLI
entailment (Section~\ref{sec:theory}A). We could not execute this
test meaningfully: with an untrained model, sampled outputs are
incoherent byte sequences, and clustering incoherent text (by any
method, NLI-based or lexical) measures surface noise, not semantic
equivalence classes. We chose not to run this and report a number
that would carry false precision; the clustering code
(\texttt{sc\_pilot.py} in the accompanying artifacts) is provided
as-is for use once a pretrained model and NLI classifier are
reachable. Third, hardware-level energy metering (Intel RAPL, NVIDIA
NVML) was not available in this compute environment (no GPU; no
access to CPU model-specific registers from within the container),
so, exactly as in Section~\ref{sec:limits}, every number in this
section remains a timing-based proxy, not a physical energy
measurement --- the proxy is simply better-justified here than in
Section~\ref{sec:methods}C, since it now contains no synthetic term.
Fourth, this is a single architecture at one size; we make no claim
that the length-mediation finding of Section~\ref{sec:pilot-length}
generalises to production-scale trained models, where learned
attention patterns, KV-cache reuse strategies, and early stopping on
end-of-sequence tokens all differ from this minimal implementation in
ways that could plausibly restore, weaken, or further explain the
semantic-content contribution to EPT. All code, the 60-prompt corpus,
and the full 540-record measurement log are included with this
submission's supplementary artifacts so that this pilot can be
rerun, extended to a pretrained model, or audited line by line.
% =============================================================================
\section{Real-Instruction-Tuned-Model Replication on Consumer GPU Hardware}
\label{sec:real}
Section~\ref{sec:pilot} closed by naming two things it could not do:
speak to a semantically competent model, and measure real hardware
energy rather than a timing-based proxy. This section does both, on
different hardware from either prior protocol, and reports what
actually happened---including that it complicates the length-mediation
finding of Section~\ref{sec:pilot-length} further than we expected,
rather than simply confirming it.
\subsection{Models, Hardware, and Backend}
\label{sec:real-setup}
We used two real, open-weight, instruction-tuned models of different
architecture and vendor: \texttt{phi3.5:3.8b} (Phi-3.5-mini-instruct,
3.8B parameters) and \texttt{gemma2:2b} (Gemma-2-2B-it, 2.6B
parameters), both served locally at Q4\_0 GGUF quantisation via
Ollama on a single consumer laptop GPU (NVIDIA RTX 3050 Laptop, 4\,GB
VRAM, 30\,W power cap; Intel Core i5-12500H; Windows 11). Direct
\texttt{transformers}-based loading with 4-bit quantisation
(\texttt{bitsandbytes}) was considered and rejected as the primary
path: it is a well-documented source of Windows/CUDA wheel failures,
and both models were already validated to run on this exact GPU via
Ollama, which we judged the lower-risk choice for a multi-hour
unattended run. Each model's own Hugging Face tokenizer
(\texttt{microsoft/Phi-3.5-mini-instruct}; for Gemma-2, the ungated
mirror \texttt{unsloth/gemma-2-2b-it}, since \texttt{google/gemma-2-2b-it}
requires authenticated licence acceptance and returns HTTP 401
otherwise---no model weights were downloaded from either repository,
only tokenizer files, used solely for token counting) supplied exact
token counts for length control (Section~\ref{sec:real-length}).
Reused unmodified from Sections~\ref{sec:methods}--\ref{sec:pilot}:
the 60-prompt corpus (\texttt{corpus.py}), the nine mutation
conditions and the SII formula (\texttt{mutation\_engine.py}), and
$\alpha=0.7$. For turnaround time, we used a stratified 15-prompt
subset (three per domain) rather than the full 60, generation output
was capped at 96 tokens (a safety ceiling, not a fixed length---see
below), and $n=255$ measurements were collected per model (15 prompts
$\times$ [1 baseline $+$ 8 non-baseline conditions $\times$ 2 length
variants]).
\subsection{Energy Measurement: Real NVML Joules, Not a Proxy}
\label{sec:real-energy}
Unlike every prior section of this paper, energy here is measured,
not modelled. A background thread polls NVIDIA's Management Library
(\texttt{pynvml.nvmlDeviceGetPowerUsage}) throughout each generation
call and the resulting power trace $P(t)$ is integrated by the
trapezoidal rule over the call's wall-clock window to obtain joules:
\begin{equation}
  E_\text{real} = \int_{t_0}^{t_1} P(t)\,dt \approx
  \sum_i \tfrac{1}{2}\big(P_i + P_{i+1}\big)\big(t_{i+1}-t_i\big).
  \label{eq:nvml}
\end{equation}
This is genuine hardware-reported power for the GPU package as a
whole---it is \emph{not} a per-process partition (any concurrent GPU
load would be included), and it does not cover CPU energy: Intel RAPL
(\texttt{/sys/class/powercap}) remains unavailable on this Windows
machine without root/Linux, exactly as in Section~\ref{sec:pilot}. No
CPU-only fallback path was invoked in this run---Ollama-plus-NVML
succeeded for all 510 generation calls across both models---but had
one been necessary it would have reused Eq.~(\ref{eq:energy})'s
TDP-proxy formula unchanged and been labelled \texttt{tdp\_proxy\_cpu}
in the output record, never silently merged with the real
\texttt{nvml\_gpu\_real} measurements. Every record additionally
carries \texttt{mean\_watts}, \texttt{max\_watts}, and the observed
sampling interval, for reviewer transparency on measurement fidelity.
\subsection{Length-Matched Control}
\label{sec:real-length}
The mutation engine has no length-control logic at all: some
mutations grow text (\texttt{noise\_verbose}, appended filler;
\texttt{ambiguity\_contradiction}, an appended clause), others
shrink or preserve it. For every non-baseline record we therefore
generated a second, length-matched variant: the mutated text
truncated (at a sentence boundary where possible) or padded with
neutral filler phrases, under the model's own tokenizer, until its
token count exactly equalled that record's own baseline token count.
Both the raw-mutated and length-matched variant were then generated
and measured independently, giving a direct, per-record causal probe
of length-mediation that is stronger than a partial correlation alone.
\subsection{Results: Both Models}
\label{sec:real-results}
\begin{table}[t]
\caption{Real-Hardware Replication: SII--EPT Correlation (Spearman),
Both Models, Raw and Length-Matched Subsets ($n=135$ each; 95\,\%
bootstrap CI in brackets where computed for the pooled column).}
\label{tab:real-pec}
\centering
\setlength{\tabcolsep}{3.2pt}
\begin{tabular}{@{}lrrr@{}}
\toprule
\textbf{Model / subset} & \textbf{PEC ($\rho$)} & \textbf{$p$} & \textbf{Partial, given tokens} \\
\midrule
Phi-3.5, raw            & $-0.050$ & $0.562$ & $-0.035$ \\
Phi-3.5, length-matched & $-0.045$ & $0.602$ & $-0.046$ \\
Phi-3.5, pooled ($n=255$) & $-0.043$ & $0.499$ & $-0.041$ \\
Gemma-2, raw            & $-0.066$ & $0.445$ & $-0.055$ \\
Gemma-2, length-matched & $+0.030$ & $0.734$ & $+0.030$ \\
Gemma-2, pooled ($n=255$) & $-0.017$ & $0.784$ & $-0.015$ \\
\bottomrule
\end{tabular}
\end{table}
Table~\ref{tab:real-pec} reports the headline result, and it is not
what either prior protocol found. On both real, semantically
competent, instruction-tuned models, the SII--EPT correlation is
statistically indistinguishable from zero---\emph{before} any length
control is applied, not only after. One-way ANOVA across the nine
mutation conditions is likewise non-significant for both models
(Phi-3.5: $F(8,246)=0.60$, $p=0.78$, $\eta^2=0.019$; Gemma-2:
$F(8,246)=1.16$, $p=0.32$, $\eta^2=0.036$; Kruskal--Wallis
corroborates both, $p=0.88$ and $p=0.55$ respectively), and Cohen's
$d$ between \texttt{baseline} and \texttt{ambiguity\_contradiction}
is negligible for both ($d=0.083$ and $d=-0.059$). This directly
contradicts the ``Contradiction Tax'' finding of
Section~\ref{sec:results} at this scale and with these models.
The length-matched control was nonetheless informative on its own
terms. A paired Wilcoxon signed-rank test between raw and
length-matched EPT for the same (prompt, mutation) pairs is
significant for Phi-3.5 ($n=120$ pairs, $W=2561$, $p=0.005$; raw EPT
higher by a mean 29.9\,mJ/token) but not for Gemma-2 ($W=3240$,
$p=0.307$; mean difference 7.2\,mJ/token). So prompt length does
measurably move EPT for at least one of the two models---just not in
a way that produces a corresponding SII--EPT correlation for either,
since the correlation was already null before length was controlled.
The SII-coefficient sensitivity analysis (Section~\ref{sec:pilot-sensitivity}'s
method, applied unmodified to this real data) confirms the null is
stable, not fragile: across the 200-draw random joint perturbation of
$\gamma$, $\delta$, $\varepsilon$, and all $b_m$ within $\pm 30\,\%$,
mean PEC stayed at $-0.032$ (SD $0.016$, range $[-0.075,\,0.006]$) for
Phi-3.5 and $-0.005$ (SD $0.033$, range $[-0.088,\,0.083]$) for
Gemma-2---no perturbation of the SII formula, for either model, pushed
the correlation anywhere near the $0.25$--$0.41$ range reported in
Sections~\ref{sec:corr} and~\ref{sec:pilot-results}.
\subsection{Semantic-Entropy Validation Pilot (Real Execution)}
\label{sec:real-sc}
Section~\ref{sec:pilot-limits} reported that \texttt{sc\_pilot.py}'s
semantic-entropy clustering test could not be meaningfully executed
against an untrained model and was left unrun. We ran it here for
real: $k=5$ samples per prompt at temperature 0.7 from each Ollama
model, clustered by bidirectional NLI entailment
(\texttt{cross-encoder/nli-deberta-v3-xsmall}, threshold 0.5,
following Kuhn et al.'s rule), on a 10-prompt $\times$ 6-condition
subset (baseline, \texttt{noise\_typo}, \texttt{ambiguity\_semantic},
\texttt{negation}, \texttt{reordering}, \texttt{ambiguity\_contradiction};
$n=60$ per model). The two models disagree: for Phi-3.5,
Spearman($\mathrm{SII},S_c$)$=-0.173$ ($p=0.187$, 95\,\% bootstrap CI
$[-0.419,\,0.084]$)---not significant. For Gemma-2,
Spearman($\mathrm{SII},S_c$)$=+0.345$ ($p=0.007$, 95\,\% CI
$[0.094,\,0.553]$)---significant and in the theoretically predicted
direction. We report both numbers as found, not averaged or
reconciled: on this small a subset, one real model shows evidence
that SII tracks actual output-level semantic entropy and the other
does not, and we do not have grounds to prefer one model's answer
over the other's as more representative of the underlying construct.
\subsection{New Limitations}
\label{sec:real-limits}
First, NVML reports whole-GPU-package power, not a per-process
partition; any other process contending for the GPU during a call
would inflate that call's measured energy, and we did not run this
replication in a controlled, single-tenant environment beyond closing
other GPU-using applications before each run. Second, Ollama serves
Q4\_0-quantised GGUF weights, not the full-precision reference
weights; any relationship (or absence of one) measured here describes
this quantised deployment configuration specifically, not an
unquantised one, and quantisation is known to change per-token
compute patterns. Third, Intel RAPL-based CPU energy accounting
remains categorically unavailable in this environment, as it was in
Section~\ref{sec:pilot}; this replication is GPU-energy-only. Fourth,
$n=255$ per model (135 raw $+$ 120 length-matched) is a modest sample
for a nine-condition comparison, and the 15-prompt stratified corpus
subset is smaller than the full 60-prompt corpus used elsewhere in
this paper, both by deliberate choice to keep total measurement time
within about an hour on this hardware. Fifth, this is two models,
not a systematic sweep across model families or scales; the
disagreement between them in Section~\ref{sec:real-sc} is itself
evidence that model-to-model variance in this relationship is large
enough that two data points cannot establish a generalisable
answer either way. Sixth, generation was capped at 96 tokens as a
safety ceiling; both models frequently produced longer answers than
this by default for these prompts, so a substantial fraction of
calls were right-censored at the cap rather than stopping naturally
at end-of-sequence---recorded per-record and available for reviewer
inspection, but a real limit on how much natural output-length
variation this replication could observe.
We take the combination of Sections~\ref{sec:pilot-length}
and~\ref{sec:real-results} to say something stronger, and less
comfortable, than either alone: the numpy pilot found the SII--EPT
correlation's \emph{magnitude} was mediated by prompt length once
measured on real (if untrained) hardware. This replication finds that
with real, semantically competent models, there may be very little
SII--EPT correlation left to mediate in the first place. The
semantic-instability construct itself may still be real and
measurable---Section~\ref{sec:real-sc}'s Gemma-2 result is some
evidence for that---without that instability reliably costing more
energy to process on the hardware and quantisation configuration
tested here. We did not adjust either finding to sit more comfortably
with the hypothesis in Section~\ref{sec:theory}B, and we do not think
a future replication should either.
% =============================================================================
\section{Discussion}
\label{sec:discussion}
\subsection{What Might Be Happening with Contradictory Prompts---and What We Cannot Yet Say}
A 79\,\% EPT premium from a single contradictory clause under the
simulation-assisted protocol (Section~\ref{sec:results}) is large
enough to invite scrutiny rather than acceptance at face value, and
Section~\ref{sec:pilot} is our attempt to provide some of that
scrutiny ourselves rather than wait for it. The gap between reordering
(18.4\,mJ/tok) and contradiction (22.2\,mJ/tok) in
Table~\ref{tab:ept} exceeds what a linear extrapolation of the SII
scale predicts. One candidate account, offered in the original
version of this analysis: reordering disrupts the \emph{form} of a
prompt while leaving content inferrable, so the model's grammatical
priors can compensate; contradiction disrupts \emph{content} while
leaving form intact, so no amount of grammatical compensation resolves
the conflict, plausibly sustaining high variance in attention-head
weight distributions across decoding steps.
Section~\ref{sec:pilot-length} complicates this account without
ruling it out. It shows that, at minimum, part of the empirical
contradiction premium is explainable by a much more mundane mechanism:
our contradiction mutation happens to lengthen the prompt more than
our reordering mutation does, and a real decoder's per-token cost
scales with context length regardless of content. We did not
instrument attention activations directly in either protocol, and an
alternative account---that contradiction mutations happen to trigger
longer effective processing in ways not fully captured by our token
and length controls---cannot be ruled out from either dataset alone.
We report the attention-variance account as a hypothesis for the
follow-up study described in Section~\ref{sec:limits}, and the
length-mediation account as a competing hypothesis this paper's own
pilot surfaced, not as a finding that resolves the question either
way.
\subsection{Threats to Validity and Limitations}
\label{sec:limits}
\textit{Construct validity of the energy proxy.} The TDP-based model
(Section~\ref{sec:methods}C) is a relative-comparison tool, not an
absolute measurement instrument, in both the simulation-assisted
protocol and the real-hardware pilot of Section~\ref{sec:pilot}.
Hardware-level metering (Intel RAPL, NVIDIA NVML) is required before
any absolute EPT figure in this paper should be treated as a physical
quantity; neither was accessible in either compute environment used
for this work. We expect relative ordering to be preserved under such
metering based on prior studies using both approaches
~\cite{samsi2023words,bannour2021evaluating}, but this is an
expectation to be tested, not a finding of this paper.
\textit{External validity of the SII base scores.} The $b_m$ values
(Table~\ref{tab:mutations}) are calibrated constants. Under the
simulation-assisted protocol, independently measured EPT tracks their
ordering exactly; under the real-hardware pilot
(Section~\ref{sec:pilot-sensitivity}), the ordering is measurably less
robust to coefficient perturbation than the aggregate PEC statistic.
They were calibrated on English-language prompts across five general
domains and may require recalibration for multilingual or
domain-specific (legal, medical, code-heavy) corpora.
\textit{Ecological validity of simulation mode.} The $\lambda$ scalar
in (\ref{eq:sim}) is more regular than live provider inference. Our
internal consistency is therefore likely higher than a live-hardware
study would show under the simulation-assisted protocol; the
real-hardware pilot in Section~\ref{sec:pilot} removes $\lambda$
entirely but substitutes an untrained model in its place, trading one
limitation for another rather than eliminating the need for further
validation.
\textit{Correlational, not causal, interpretation.} This is the
threat we want to be most direct about. H1 and H2 are tested as
statistical associations under two different controlled protocols, on
two different single-hardware configurations, neither of which
manipulates model internals directly. Section~\ref{sec:pilot-length}'s
length-mediation finding is itself evidence that at least one
non-semantic confound can materially affect this kind of measurement,
which is precisely why we do not claim to have demonstrated the
Landauer-motivated mechanism in Section~\ref{sec:theory}B---only that
the data across both protocols are directionally consistent with it,
while also being consistent with a more mundane length-driven
explanation that has not yet been ruled out. The necessary next step,
prioritised over all others, is replication with RAPL/NVML hardware
metering on a \emph{pretrained} open-weight model, with mutation
conditions length-matched by padding or truncation, so that the
length-mediation and semantic-instability contributions to EPT can be
separated rather than confounded.
\subsection{Practical Implications, Stated Conditionally}
Pending the hardware- and pretrained-model validation above, three
implications follow \emph{if} the observed associations hold under
stricter measurement. \textbf{Contradiction screening} may be worth
the cost of automated NLI pre-submission checks
~\cite{farquhar2024detecting} at large query volumes, given the size
of the observed premium under both protocols, though
Section~\ref{sec:pilot-length} suggests some of that premium may
reduce to a simpler recommendation: keep contradictory or corrective
prompt edits concise rather than appended. \textbf{Verbose prompts}
appear to cost more than token count alone would predict under the
simulation-assisted protocol; the real-hardware pilot is consistent
with verbosity's cost operating substantially through added length,
which still supports a conciseness-review recommendation, if for a
more mundane reason than originally proposed. \textbf{SII logging} is
cheap enough (Section~\ref{sec:methods}B) to run in production
regardless of whether the underlying mechanism is ultimately
semantic or length-driven, since it functions as a leading indicator
for prompt-quality drift either way.
% =============================================================================
\section{Conclusion}
Under a controlled, proxy-based protocol, prompt semantic instability
is statistically associated with per-token inference energy. The
Semantic Instability Index quantifies prompt difficulty in a single
pass; the Prompt Entropy Coefficient places the SII--EPT relationship
on a rank-correlation footing with bootstrap confidence intervals and
an a priori power analysis; the Mutation Engine provides a
reproducible nine-class protocol for generating semantically graded
prompt variants. Across 540 measurements under the simulation-assisted
protocol, both pre-specified hypotheses are supported: PEC $=0.408$
($p<0.001$) for H1, and $\eta^2=0.157$ for H2, with logically
contradictory prompts showing a 79.0\,\% EPT premium not explained by
output token count. An independently executed, real-hardware pilot
replication (Section~\ref{sec:pilot}), built from scratch to avoid any
synthetic energy term, reproduces the direction of both effects at
smaller magnitude ($\mathrm{PEC}=0.252$, $\eta^2=0.086$) and surfaces a
specific, important qualification: much of the association is
mediated by mutation-induced prompt-length changes rather than by
semantic content once length is controlled for. A coefficient
sensitivity analysis on that same real data shows the aggregate
SII--EPT direction is robust to $\pm 30\,\%$ perturbation of every
tunable constant, while the specific mutation-by-mutation ordering is
not. We present these first two protocols as findings that motivate,
rather than conclude, a research programme: on their own, they say
the question is real, measurable, and only partially resolved.
The follow-on work prioritised above---real hardware metering on a
pretrained model with explicit length control, and a real execution
of the semantic-entropy validation pilot---is exactly what
Section~\ref{sec:real} then carried out, on two open-weight
instruction-tuned models (Phi-3.5-mini, Gemma-2-2B) with real
NVML-measured GPU energy. It does not confirm the earlier findings at
smaller scale; it substantially weakens them. The SII--EPT correlation
is statistically indistinguishable from zero for both models, both
before and after length matching ($|\rho|\le 0.07$, all $p>0.4$), and
the nine-condition ANOVA is non-significant for both. The
semantic-entropy pilot is mixed rather than confirmatory: SII tracks
measured output-level semantic entropy significantly for Gemma-2
($\rho=0.345$, $p=0.007$) but not for Phi-3.5 ($\rho=-0.173$,
$p=0.187$). Taken together across all three protocols, we think the
honest summary is this: a proxy-measured, simulation-assisted setup
and an untrained real-hardware transformer both show a positive
SII--EPT association that is partly a length artifact; two real,
semantically competent models, measured with real GPU energy, show
essentially no SII--EPT association at all, length-matched or not.
The semantic-instability construct may still be measurable at the
output level (Section~\ref{sec:real-sc}), but on the evidence gathered
across all three protocols in this paper, we no longer think prompt
semantic instability can be presented as a demonstrated correlate of
per-token inference energy in real, deployed models. That is a
materially different conclusion from the one this paper's abstract
would have drawn after only the first two protocols, and we think
reporting that shift plainly is more useful to the field than
retaining the earlier framing would have been.
Remaining follow-on work: (1) RAPL-based CPU energy accounting, still
unavailable on the Windows environment used for Section~\ref{sec:real}
and requiring a Linux host with root access; (2) a systematic sweep
across more model families and scales than the two tested here, given
how much the two disagreed in Section~\ref{sec:real-sc}; (3) a larger
corpus and measurement count than the 15-prompt, $n=255$-per-model
scale used in Section~\ref{sec:real}, to obtain tighter confidence
intervals around what are, at this scale, statistically null results
rather than precisely estimated small effects; and (4) extension to
multi-turn conversation settings, where context entropy may
accumulate across turns and interact with KV-cache reuse in ways none
of the three protocols in this paper capture.
% =============================================================================
\balance
\begin{thebibliography}{00}
\bibitem{strubell2019energy}
E.~Strubell, A.~Ganesh, and A.~McCallum,
``Energy and policy considerations for deep learning in NLP,''
in \textit{Proc.\ 57th Annu.\ Meet.\ Assoc.\ Comput.\ Linguist.\ (ACL)},
Florence, Italy, Jul.~2019, pp.~3645--3650.
DOI:~10.18653/v1/P19-1355.
\bibitem{patterson2022carbon}
D.~Patterson et al.,
``Carbon emissions and large neural network training,''
\textit{Commun.\ ACM}, vol.~65, no.~6, pp.~35--35, Jun.~2022.
DOI:~10.1145/3522589.
\bibitem{luccioni2023estimating}
A.~S.~Luccioni, S.~Viguier, and A.-L.~Ligozat,
``Estimating the carbon footprint of BLOOM, a 176B parameter
language model,''
\textit{J.~Mach.\ Learn.\ Res.}, vol.~24, no.~253,
pp.~1--15, 2023.
arXiv:2211.02001.
\bibitem{wu2022sustainable}
C.-J.~Wu et al.,
``Sustainable AI: Environmental implications, challenges,
and opportunities,''
in \textit{Proc.\ 5th Conf.\ Mach.\ Learn.\ Syst.\ (MLSys)},
Santa Clara, CA, USA, 2022.
arXiv:2111.00364.
\bibitem{samsi2023words}
S.~Samsi et al.,
``From words to watts: Benchmarking the energy costs of large
language model inference,''
in \textit{Proc.\ IEEE High Perform.\ Extreme Comput.\ Conf.\
(HPEC)}, Dedham, MA, USA, Sep.~2023, pp.~1--9.
DOI:~10.1109/HPEC58863.2023.10363447.
\bibitem{bannour2021evaluating}
N.~Bannour, S.~Ghannay, A.~N\'ev\'eol, and A.-L.~Ligozat,
``Evaluating the carbon footprint of NLP methods: A survey
and analysis of existing tools,''
in \textit{Proc.\ 2nd Workshop SustaiNLP (EMNLP)},
Nov.~2021, pp.~11--21.
DOI:~10.18653/v1/2021.sustainlp-1.2.
\bibitem{desislavov2023trends}
R.~Desislavov, F.~Mart\'inez-Plumed, and J.~Hern\'andez-Orallo,
``Trends in AI inference energy consumption: Beyond the
performance-vs-parameter laws of deep learning,''
\textit{Sustain.\ Comput.: Inform.\ Syst.}, vol.~38,
art.~100857, Apr.~2023.
DOI:~10.1016/j.suscom.2023.100857.
\bibitem{faiz2024llmcarbon}
A.~Faiz, S.~Kaneda, R.~Wang, R.~Osi, P.~Sharma, F.~Chen, and
L.~Jiang,
``LLMCarbon: Modeling the end-to-end carbon footprint of large
language models,''
in \textit{Proc.\ 12th Int.\ Conf.\ Learn.\ Represent.\ (ICLR)},
Vienna, Austria, May~2024.
arXiv:2309.14393.
\bibitem{chien2023reducing}
A.~A.~Chien, L.~Lin, H.~Nguyen, V.~Rao, T.~Sharma, and
R.~Wijayawardana,
``Reducing the carbon impact of generative AI inference (today
and in 2035),''
in \textit{Proc.\ 2nd Workshop Sustain.\ Comput.\ Syst.\
(HotCarbon)}, Boston, MA, USA, Jul.~2023.
DOI:~10.1145/3604930.3605705.
\bibitem{sahoo2024systematic}
P.~Sahoo, A.~K.~Singh, S.~Saha, V.~Jain, S.~Mondal, and
A.~Chadha,
``A systematic survey of prompt engineering in large language
models: Techniques and applications,''
\textit{arXiv preprint}, 2024.
arXiv:2402.07927.
\bibitem{landauer1961irrev}
R.~Landauer,
``Irreversibility and heat generation in the computing process,''
\textit{IBM J.~Res.\ Dev.}, vol.~5, no.~3,
pp.~183--191, Jul.~1961.
DOI:~10.1147/rd.53.0183.
\bibitem{berut2012experimental}
A.~B\'erut, A.~Arakelyan, A.~Petrosyan, S.~Ciliberto,
R.~Dillenschneider, and E.~Lutz,
``Experimental verification of Landauer's principle linking
information and thermodynamics,''
\textit{Nature}, vol.~483, no.~7388, pp.~187--189, Mar.~2012.
DOI:~10.1038/nature10872.
\bibitem{bennett1982thermo}
C.~H.~Bennett,
``The thermodynamics of computation---a review,''
\textit{Int.~J.~Theor.\ Phys.}, vol.~21, nos.~12,
pp.~905--940, 1982.
DOI:~10.1007/BF02084158.
\bibitem{vaswani2017attention}
A.~Vaswani et al.,
``Attention is all you need,''
in \textit{Adv.\ Neural Inf.\ Process.\ Syst.\ (NeurIPS)},
vol.~30, Long Beach, CA, USA, Dec.~2017, pp.~5998--6008.
arXiv:1706.03762.
\bibitem{kuhn2023semantic}
L.~Kuhn, Y.~Gal, and S.~Farquhar,
``Semantic uncertainty: Linguistic invariances for uncertainty
estimation in natural language generation,''
in \textit{Proc.\ 11th Int.\ Conf.\ Learn.\ Represent.\ (ICLR)},
Kigali, Rwanda, May~2023.
arXiv:2302.09664.
\bibitem{farquhar2024detecting}
S.~Farquhar, J.~Kossen, L.~Kuhn, and Y.~Gal,
``Detecting hallucinations in large language models using
semantic entropy,''
\textit{Nature}, vol.~630, no.~8017, pp.~625--630, Jun.~2024.
DOI:~10.1038/s41586-024-07421-0.
\bibitem{malinin2021uncertainty}
A.~Malinin and M.~Gales,
``Uncertainty estimation in autoregressive structured
prediction,''
in \textit{Proc.\ 9th Int.\ Conf.\ Learn.\ Represent.\ (ICLR)},
May~2021.
arXiv:2002.07650.
\bibitem{schwartz2020green}
R.~Schwartz, J.~Dodge, N.~A.~Smith, and O.~Etzioni,
``Green AI,''
\textit{Commun.\ ACM}, vol.~63, no.~12, pp.~54--63, Dec.~2020.
DOI:~10.1145/3381831.
\bibitem{dodge2022measuring}
J.~Dodge et al.,
``Measuring the carbon intensity of AI in cloud instances,''
in \textit{Proc.\ 2022 ACM Conf.\ Fairness, Accountability,
Transparency (FAccT)}, Seoul, Korea, Jun.~2022,
pp.~1877--1894.
DOI:~10.1145/3531146.3533234.
\bibitem{flesch1948readability}
R.~Flesch,
``A new readability yardstick,''
\textit{J.~Appl.\ Psychol.}, vol.~32, no.~3, pp.~221--233,
1948.
DOI:~10.1037/h0057532.
\bibitem{lacoste2019quantifying}
A.~Lacoste, A.~Luccioni, V.~Schmidt, and T.~Dandres,
``Quantifying the carbon emissions of machine learning,''
in \textit{Workshop Tackling Climate Change with AI (NeurIPS)},
Vancouver, Canada, Dec.~2019.
arXiv:1910.09700.
\bibitem{cohen1988statistical}
J.~Cohen,
\textit{Statistical Power Analysis for the Behavioral Sciences},
2nd~ed. Hillsdale, NJ, USA: Lawrence Erlbaum Associates, 1988.
\end{thebibliography}
\end{document}
